\begin{abstract}

Zero-shot recognition models require extensive training data for generalization.
However, in zero-shot 3D classification, collecting 3D data and captions is costly and labor-intensive, posing a significant barrier compared to 2D vision.
% To fix this issue, recent research on zero-shot 3D classification focuses on distilling knowledge from a pretrained visual language model, but this isn’t a complete solution to the data scarcity.
Recent advances in generative models have achieved unprecedented realism in synthetic data production, and recent research shows the potential for using generated data as training data.
% In contrast, advances in generative models now enable creating realistic synthetic data, which could be utilized for training.
Here, naturally raising the question: Can synthetic 3D data generated by generative models be used as expanding limited 3D datasets? 
In response, we present a synthetic 3D dataset expansion method, \textbf{Te}xt-guided \textbf{G}eometric \textbf{A}ugmentation (\textbf{TeGA}).
TeGA is tailored for language-image-3D pretraining, which achieves SoTA in zero-shot 3D classification, and uses a generative text-to-3D model to enhance and extend limited 3D datasets.
Specifically, we automatically generate text-guided synthetic 3D data and introduce a consistency filtering strategy to discard noisy samples where semantics and geometric shapes do not match with text. 
In the experiment to double the original dataset size using TeGA, our approach demonstrates improvements over the baselines, achieving zero-shot performance gains of 3.0\% on Objaverse-LVIS, 4.6\% on ScanObjectNN, and 8.7\% on ModelNet40.
These results demonstrate that TeGA effectively bridges the 3D data gap, enabling robust zero-shot 3D classification even with limited real training data and paving the way for zero-shot 3D vision applications. 

\end{abstract}


\begin{comment}
While zero-shot recognition models for open-vocabulary 3D understanding require extensive training data for generalization, the collection of both 3D data and their captions remains costly and labor-intensive, creating a significant barrier compared to the data abundance in 2D vision. 
Simultaneously, recent advances in generative models have achieved unprecedented realism in synthetic data production, naturally raising the question: To what extent can synthetic data serve as an effective substitute for real 3D data? % in 3D classification? 
In response, we present a synthetic 3D dataset expansion method, \textbf{Te}xt-guided \textbf{G}eometric \textbf{A}ugmentation (\textbf{TeGA}), designed to use generative text-to-3D models to expand limited 3D datasets and apply text guidance for multimodal zero-shot 3D classification. 
Specifically, we automatically generate text-guided synthetic 3D data %from an off-the-shelf text-to-3D model 
and introduce a consistency filtering strategy to discard noisy samples where semantics and geometric shape do not match. \tbd{Surprisingly, TeGA successfully enhances dataset expansion: inserting synthetic 3D data generated by TeGA for ModelNet40 with $<40$ original instances matches the use of 100\% real data. Our approach demonstrates improvements over the baselines, achieving zero-shot performance gains of 2.2\% on Objaverse-LVIS, 2.2\% on ScanObjectNN, and 8.8\% on ModelNet40.} These results demonstrate that TeGA effectively bridges the 3D data gap, enabling robust zero-shot 3D classification even with limited real training data and paving the way for scalable 3D vision applications. 
\end{comment}